\documentclass[../Main.tex]{subfiles}

\begin{document}

\section{Evaluation} \label{sec:Evaluation}

In diesem Kapitel wird überprüft, ob die acht Designvorschläge (V1--V8), die wir aus unseren Spieleanalysen und der Leveldesign-Recherche hergeleitet und in Kapitel~\ref{sec:Umsetzung} umgesetzt haben, die beabsichtigten Wirkungen auf Spielfluss, Schwierigkeit und Spielgefühl tatsächlich erzeugen.

Wir sind beim Playtest wie in Abschnitt~\ref{sec:VorgehenPlaytest} beschrieben vorgegangen: Zuerst wurde die Testperson mit dem Pre-Questionnaire eingeschätzt, danach folgte der Playtest, der aufgenommen wurde, und anschliessend das Post-Questionnaire sowie das semistrukturierte Interview für tiefergehende Fragen.

Das folgende Kapitel beschreibt zunächst den Ablauf der Playtests. Danach gehen wir die Designvorschläge einzeln durch, stellen jeweils die Ergebnisse dar und interpretieren sie. Anschliessend ordnen wir weitere für die Erkenntnisfrage relevante Beobachtungen sowie den Gesamteindruck ein. Die Verbesserungen, die wir aus den Playtests ableiten, folgen in Abschnitt~\ref{sec:Verbesserungen}.

\subsection{Durchführung und Testpersonen} \label{sec:Durchfuehrung}

Alle fünf Playtests wurden von derselben Person geleitet, um möglichst vergleichbare Ergebnisse zu erhalten. Gespielt wurde immer mit Maus und Tastatur. Zwei Playtests fanden über Discord mit Bildschirmübertragung statt, drei vor Ort. Alle Testpersonen spielten dieselbe exportierte Version des Spiels (für Mac oder Windows) und benutzten dabei ihre gewohnte Hardware. Aufgenommen haben wir die Playtests jeweils mit einem Handy.

Die Playtests fanden zwischen dem 26.10.2026 und dem 02.11.2026 statt. Das Pre-Questionnaire dauerte in der Regel höchstens fünf Minuten. Der Playtest selbst war unterschiedlich lang: Der kürzeste dauerte etwa 20, der längste etwa 40 Minuten. Das lag unserer Einschätzung nach weniger am Können der Testpersonen als an ihren Absichten: Manche wollten Glitches finden, andere das Spiel möglichst schnell durchspielen. Post-Questionnaire und Interview dauerten zusammen etwa 15 Minuten. Insgesamt nahm jeder Playtest für die Testperson also etwa 40 bis 60 Minuten in Anspruch.

Während des Spielens haben wir keinen Smalltalk geführt, um die Testpersonen nicht abzulenken. In entspannten Phasen haben wir mit ihnen allerdings über das Spiel gesprochen. Hatte eine Testperson eine Frage zum Kampfsystem oder etwas vergessen, haben wir die Frage beantwortet. Ganz zu Beginn haben wir den Testpersonen zudem gezeigt, welche Aktionen es gibt und wie man sie ausführt, da der Prototyp kein Tutorial enthält.

Obwohl wir die Testpersonen nicht zum lauten Denken aufgefordert hatten (Abschnitt~\ref{sec:MethodePlaytest}), haben sie während des Spielens immer wieder von sich aus erzählt, was sie gerade denken. Dadurch erhielten wir schon in der Beobachtung viele Informationen über ihre Gedanken und Absichten. Das semistrukturierte Interview fiel deshalb deutlich kürzer aus als geplant, da meist nur wenige Fragen offen blieben.

Nach der Auswertung haben wir die Aufnahmen gelöscht (Abschnitt~\ref{sec:AuswahlTP}) und die Testpersonen in dieser Arbeit anonymisiert. So lassen sich die einzelnen Aussagen nicht direkt mit einer Person in Verbindung bringen. Für die Auswertung spielt es zudem keine Rolle, wer eine Aussage gemacht hat, da wir die Testpersonen über das Pre-Questionnaire charakterisiert haben. Tabelle~\ref{tab:Testpersonen} fasst diese Charakterisierung zusammen; sie stützt sich auf die Antworten P1--P8 und wird durch unsere Beobachtungen ergänzt (Anhang~\ref{sec:PlaytestProtokolle}).\footnote{Mit Claude formuliert, siehe KI-Deklaration (Abschnitt~\ref{sec:KIDeklaration}).}

\begin{table}[H]
    \centering
    \small
    \renewcommand{\arraystretch}{1.3}
    \begin{tabular}{|>{\raggedright\arraybackslash}p{1.2cm}|>{\raggedright\arraybackslash}p{12.6cm}|}
        \hline
        \textbf{TP} & \textbf{Charakterisierung} \\
        \hline
        TP1 & Spielt während der Schulzeit kaum, in den Ferien dafür viel, vor allem Survival- und Shooter-Spiele, und schätzt das eigene Können mit 3--4 ein. Bevorzugt aggressives Spiel, zog sich im Playtest bei viel Schaden aber zurück, heilte regelmässig und suchte nebenbei nach Glitches. \\
        \hline
        TP2 & Spielt unter der Woche etwa 3--4 Stunden, in den Ferien deutlich mehr, und zwar quer durch viele Genres (u.\,a. Souls-likes, Shooter, Horror) bei einer Selbsteinschätzung von 3--4. Spielte wie angegeben sehr aggressiv und nahm viele Treffer in Kauf, ging in kleinen Räumen aber taktischer vor. \\
        \hline
        TP3 & Spielt je nach Neuerscheinungen 0 bis 20 Stunden pro Woche, vor allem Action- und Platformer-Spiele, und schätzt sich mit 3 ein. Spielt nach eigener Aussage ohne viel nachzudenken; im Playtest sehr aggressiv, die eigenen Lebenspunkte wurden oft ignoriert. \\
        \hline
        TP4 & Spielt heute kaum noch, früher mehrere Stunden pro Woche, am liebsten Taktik- und Aufbauspiele, und schätzt sich mit 3 ein. Mag je nach Stimmung beide Spielstile, spielte im Playtest aber eher wenig taktisch und kommentierte das Spiel ausführlich. \\
        \hline
        TP5 & Spielt 4 bis 30 Stunden pro Woche in vielen Genres, darunter Souls-likes und Roguelikes, und schätzt sich mit 4 am höchsten ein. Spielt methodisch und probiert alles aus; im Playtest als einzige Testperson eher taktisch, mit viel Interesse an Mustern und Glitches. \\
        \hline
    \end{tabular}
    \caption{Charakterisierung der Testpersonen (Pre-Questionnaire P1--P8, ergänzt durch Beobachtungen)}
    \label{tab:Testpersonen}
\end{table}

\subsection{Ergebnisse und Diskussion pro Designvorschlag} \label{sec:ErgVorschlag} \label{sec:Interpretation}

Im Folgenden gehen wir die Designvorschläge V1--V8 der Reihe nach durch. Zu jedem Vorschlag beschreiben wir zuerst kurz das Ziel, wie wir es schon in Kapitel~\ref{sec:Umsetzung} hergeleitet haben. Danach führen wir als Ergebnisse die Antworten und Beobachtungen auf, die für diesen Vorschlag wichtig sind. Zum Schluss interpretieren wir die Ergebnisse: Wir halten zuerst fest, ob der Vorschlag bestätigt, teilweise bestätigt oder nicht bestätigt wurde, und erklären danach, warum. Die einzelnen Antworten der Testpersonen sind im Anhang aufgeführt (Abschnitt~\ref{sec:PlaytestProtokolle}). Verweise wie (TP3, B1) beziehen sich auf diese Protokolle: P steht für das Pre-Questionnaire, Q für das Post-Questionnaire, B für das Beobachtungsraster und I für das Interview.

\begin{itemize}[leftmargin=1.5cm]
    \item \textcolor{red}{Verweise auf den Anhang im Format (TP3, B1) bei allen Vorschlägen ergänzen, so wie bei V1.}
\end{itemize}

\subsubsection{V1: Überlegenheit als Grundzustand}

\paragraph{Ziel}
Der Spieler soll sich seinen Gegnern grundsätzlich überlegen fühlen, wobei dieser Zustand gezielt durchbrochen werden kann. Aggression soll sich dadurch sicher und belohnt anfühlen, während Spannung erst dort entsteht, wo die Kontrolle kurzzeitig bedroht wird (Abschnitt~\ref{sec:V1}).

\paragraph{Ergebnisse}
Drei von fünf Testpersonen stimmten der Aussage, sich den Gegnern grundsätzlich überlegen gefühlt zu haben, völlig zu (TP2, TP3 und TP4, jeweils Q1); TP1 und TP5 bewerteten sie neutral. Zugestimmt haben vor allem jene Testpersonen, die im Playtest sehr aggressiv spielten. TP3 ignorierte beispielsweise ihre eigenen Lebenspunkte und konzentrierte sich nur auf die Gegner (TP3, B1). Am besten fand TP3 unter anderem, dass man einen Archer mit einem schweren Angriff sofort besiegen kann (TP3, Q12). TP4 meinte: \enquote{Ich geniesse meine Überlegenheit}; besonders befriedigend sei es, alle Gegner auf einmal zu besiegen (TP4, B11), und Horden von Gegnern zu zerstören, habe sich am besten angefühlt (TP4, Q12). TP5 spielte dagegen weniger aggressiv, eher taktisch und methodisch (TP5, B1), und hätte sich den Spieler im Vergleich zu den Gegnern noch stärker vorgestellt, etwa wie in \game{Hyrule Warriors} (TP5, Q14). TP1 bevorzugt laut eigener Angabe zwar aggressives Spiel (TP1, P6), sprang im Playtest aber aus dem Kampf heraus, sobald TP1 zu viel Schaden nahm (TP1, B1). Auch TP5 zog sich zurück, wenn sie Schaden nahm (TP5, B4). Wie TP5 beobachtete TP1 das Spiel viel und suchte nach Glitches (TP1, B13; TP5, B13).

Die Überlegenheit wurde aber auch immer wieder durchbrochen. Insbesondere wenn die Testpersonen von einer Gegnergruppe umzingelt waren, konnte es passieren, dass sie wie in einem Stunlock gefangen waren und erst nach erlittenem Schaden wieder herauskamen (TP4, B2; TP5, Q13). In der Luft waren sie zudem durch die Archers verwundbar: Wurden sie getroffen, fielen sie zu Boden, und es kam vor, dass sie beim Aufstehen weiteren Schaden erlitten (TP2, I6; TP4, B2 und B11; TP5, I8). Zudem empfand jede Testperson einen bestimmten Archer-Typ als besonders unangenehm, aber nicht alle denselben: TP2, TP3 und TP4 den Heavy Archer (TP2, B12; TP3, B12; TP4, B11), TP1 den Chain Archer (TP1, I1b) und TP5 die Archers allgemein (TP5, B3). Auch die Trägheit der Combos machte die Testpersonen anfällig (TP3, I6; TP4, Q13). Während einer Combo spielte dies eine geringere Rolle: Die Spielfigur schlägt dabei nicht nur geradeaus, sondern um sich, sodass kein Nahkampfgegner aus irgendeiner Richtung herankommen und einen Treffer landen konnte. Am Ende einer Combo konnte es jedoch passieren, dass man Schaden erlitt. Die Archers konnten den Spieler dagegen auch während einer Combo treffen: Kam ein Pfeil angeflogen, konnte man fast nicht ausweichen, und vier von fünf Testpersonen haben den Schaden einfach in Kauf genommen (TP1, TP2, TP3 und TP4, jeweils B4).

Der Heavy Soldier, der die Überlegenheit gezielt durchbrechen sollte, wurde dagegen kaum als Bedrohung erlebt. Beim ersten Auftreten wirkte er zwar eindrücklich: TP1 sah schockiert aus (TP1, B9), und TP2 war vom ersten Treffer überrascht (TP2, B9). Danach fanden ihn die Testpersonen aber einfacher als erwartet (TP1, I1a) und bemerkten, dass er häufig verfehlt (TP1, B9; TP3, I5; TP4, I2). Sie griffen ihn ohne zu zögern an oder hoben ihn sich für den Schluss auf (TP2, I5; TP3, B3; TP5, B9).

Angespannt waren die Testpersonen vor allem dann, wenn ihre Kontrolle bedroht war. TP3 geriet nur in Panik, wenn sie wenig Leben hatte und gleichzeitig viele Gegner da waren (TP3, B2); solange noch Fernkämpfer im Raum waren, empfand sie den Kampf als hektisch, danach als entspannt (TP3, I7). Für TP5 hing die Anspannung vor allem an den eigenen Lebenspunkten: \enquote{Wenig Leben = Stress} (TP5, I7). Gleichzeitig meinte TP5, nie unter starkem Stress gewesen zu sein (TP5, B11). TP1 sagte in Raum~11 voraus, gleich zu sterben (TP1, B11). TP4 war nach eigener Aussage dagegen durchgehend \enquote{locked in} (TP4, I7).

\paragraph{Interpretation}
Insgesamt sehen wir V1 als teilweise bestätigt: Der Grundzustand der Überlegenheit hat sich bestätigt, das gezielte Durchbrechen dieses Zustands nur teilweise.

Im Grundzustand konnten wir diesen Vorschlag bestätigen: Bei den Testpersonen, die aggressiv spielten, trifft die Aussage völlig zu; sie genossen ihre Überlegenheit. Für TP5, die eher taktisch spielte, war der Spieler dagegen noch etwas zu schwach und hätte stärker sein können. TP1 schätzte sich zwar als eher aggressiv ein, spielte nach TP5 aber am taktischsten und bewertete die Überlegenheit wie TP5 neutral. Dafür sehen wir zwei mögliche Gründe: Entweder reichten unsere Massnahmen nicht aus, um TP1 zu aggressivem Spiel zu bewegen, oder TP1 war davon abgelenkt, dass es sich um einen Prototyp handelt, suchte deshalb eher nach Glitches und war so in einer taktischeren Denkweise. Für diese beiden eher taktisch spielenden Testpersonen fühlte sich Aggression nicht sicher an; beide zogen sich bei Schaden zurück. Die einzelnen Gegner waren wie geplant nie ein Problem, sondern nur in Kombination und in grossen Mengen.

Die meisten der geplanten Brüche (Abschnitt~\ref{sec:V1}) haben gewirkt. Gruppen waren gefährlicher als einzelne Gegner, auch weil wir den Gegnermanager sehr aggressiv eingestellt haben, und einige Testpersonen wichen ihnen aus oder mussten sich zwischendurch erholen (Abschnitt~\ref{sec:ErgVorschlag}, V2). Das Heilen brachte einen taktischen Bruch in den eher aggressiv ausgelegten Kampf, weil man sich dafür einen sicheren Platz suchen musste. Dass für manche Gegner mehrere Combos nötig waren, zwang den Spieler ebenfalls, etwas taktischer zu denken; man konnte aber trotzdem einfach drauflosschlagen. Bei manchen Testpersonen hatte das Erfolg, andere zogen sich zurück.

Auch die Trägheit der Combos brach die Überlegenheit. Hier unterscheiden wir aber von den anderen Brüchen: Die Trägheit hat mit Abstand die meisten Kommentare ausgelöst, war aber kein taktischer Bruch, bei dem die Testpersonen meinten, sie hätten anders handeln können. Sie garantierte vielmehr, getroffen zu werden, denn auch wer ausweichen wollte, steckte noch in der Animation fest. Beim Heilen ist es ähnlich, in den meisten Räumen gab es jedoch Boxen, auf denen man es taktisch einsetzen konnte. Das Heilen liegt damit auf der Grenze zwischen beiden Arten von Brüchen.

Nicht gelungen ist der Bruch durch den Heavy Soldier, denn er ist etwas zu schwach geraten. Wir haben ihn als stärkere Variante des Soldiers entworfen, daher auch der Name, und ihm mehr Lebenspunkte und mehr Schaden gegeben; zudem lässt er sich als einziger Gegner nicht unterbrechen (Abschnitt~\ref{sec:Gegnertypen}). Um ihn etwas abzuschwächen, greift er aber langsamer an und trifft ungenauer, sodass er häufiger verfehlt als der Soldier. Diese Abschwächungen und der Umstand, dass wir Soldiers meist in Horden, Heavy Soldiers aber nur selten in grösserer Zahl erscheinen liessen, machten einen Heavy Soldier leichter als eine Horde normaler Soldiers. Man konnte ihm leicht davonrennen und ihn erst zum Schluss besiegen. Dem hätten wir entgegenwirken können, indem wir ihn wie den Soldier häufiger und in Gruppen platziert hätten; das ist uns vor dem Playtest aber nicht aufgefallen. Als Bruch wirkte der Heavy Soldier damit nicht. Gelungen ist dagegen seine Inszenierung: Weil wir ihn in einem kleinen Raum eingeführt haben, wirkte er beim ersten Auftreten stärker, als er tatsächlich ist.

Dafür brachen Faktoren die Überlegenheit, die wir nicht als Bruch geplant hatten: die Treffer in der Luft und die Archers. Dass jede Testperson einen anderen Archer-Typ als besonders unangenehm empfand, liegt wahrscheinlich an ihren Bewegungsmustern, an der Reihenfolge, in der sie die Gegner besiegten, oder an den Stellen, an denen sie zu kämpfen begannen. Spannung entstand damit zwar dort, wo die Kontrolle bedroht war, vor allem bei wenig Leben und vielen Gegnern, aber zu einem guten Teil durch Faktoren, die wir nicht gezielt gesetzt hatten.

Den Kontrollzustand erlebten die Testpersonen damit ähnlich wie in \game{NieR: Automata}, also als Kontrolle im Grundzustand (Abschnitt~\ref{sec:Vergleich}).

\begin{itemize}[leftmargin=1.5cm]
    \item \textcolor{red}{\textbf{Noch offen bei V1:}}
    \item \textcolor{red}{\textbf{Gestrichen:} \enquote{Dass diese Faktoren die Überlegenheit brechen, hatten wir bereits vor dem Playtest bemerkt, sie aber als schwächer und nur ergänzend eingeschätzt.} Der Satz bezog sich auf Trägheit, Luft und Gruppen; Trägheit und Gruppen sind jetzt geplante Brüche (Abschnitt~\ref{sec:V1}). Falls ihr das für die Treffer in der Luft so gemeint habt, dort ergänzen.}
    \item \textcolor{red}{\textbf{Heavy Soldier:} Euer Satz \enquote{das hat auch TP5 angemerkt, und alle Testpersonen teilten diese Meinung} ist im Anhang nicht belegt (TP5 sagt nur, er springe irgendwohin, B11). Belegt ist der starke erste Eindruck bei TP1 und TP2 (B9). Ich habe den Satz deshalb weggelassen; falls aus der Aufnahme, im Anhang ergänzen.}
\end{itemize}

\subsubsection{V2: Begrenzung gleichzeitig aktiver Gegner}

\paragraph{Ziel}
In Encountern mit vielen Gegnern sollen nicht alle Gegner gleichzeitig erscheinen oder gleichzeitig angreifen. So sollen auch grosse Gruppen lesbar und fair bleiben, statt den Spieler zu überfordern (Abschnitt~\ref{sec:V2}).

\paragraph{Ergebnisse}
Vier von fünf Testpersonen stimmten der Aussage, auch bei vielen Gegnern die Kontrolle behalten zu haben, eher oder völlig zu (TP2 bis TP5, jeweils Q2); TP1 stimmte eher nicht zu (TP1, Q2). Sie hatten dabei alle eine eigene Strategie, um mit grösseren Gegnermassen umzugehen. Einerseits liessen sich ähnliche Muster beobachten: Einige griffen einfach den ersten Gegner an, der bei ihnen war, andere wollten zuerst alle Archers loswerden. Andererseits gab es unterschiedliche Ausprägungen: TP2 rannte weg, um sich heilen zu können, während TP5 auf die Boxen sprang, die sie gleich zu Beginn eines Encounters von Gegnern befreit hatte, um einen sicheren Rückzugsort zu haben. TP3 kämpfte immer gegen viele Gegner gleichzeitig; dabei kam es vor, dass die Testperson einen kurzen Moment der Panik hatte, wenn sie zu wenig Leben hatte. TP5 versuchte eher, Gruppen zu meiden und nur einzelne Gegner anzugreifen.

TP2 meinte, es sei schwerer, wenn die Gegner in einem Haufen seien, als wenn sie einzeln kämen, aber nicht unangenehm schwer. TP5 meinte, sie könne nicht einschätzen, wie viel Schaden eine Gruppe von Gegnern mache, wenn sie in die Gruppe hineinlaufe, und ob sie dann im Haufen auch noch in einen Stunlock gerate \textcolor{red}{[Satz endete mit \enquote{und deswegen} -- vervollständigen]}.

\begin{itemize}[leftmargin=1.5cm]
    \item \textcolor{red}{Korrigiert: In eurem Text stand zweimal \enquote{TP2}. Das Säubern der Boxen als Rückzugsort war laut Anhang TP5 (B3), bitte prüfen.}
    \item \textcolor{red}{Fehlt: TP4 häufte Soldiers beim Archer an, um alle auf einmal zu treffen (B2).}
    \item \textcolor{red}{Der Satz zu den Mustern (erster Gegner / zuerst Archers) gehört eher zu V3; hier ggf. kürzen.}
\end{itemize}

\begin{itemize}[leftmargin=1.5cm]
    \item \textcolor{red}{\textbf{Neu -- Ergebnisse (Daten aller fünf Testpersonen):}}
    \item \textcolor{red}{\textbf{TP1 ergänzen:} rannte hin und her, um eine Welle zu zerstreuen, und rannte weg, wenn unklar war, wie viele Gegner da sind (B2); wurde von Heavy Soldiers in eine Ecke gedrängt (B6).}
    \item \textcolor{red}{Fehlt noch (rote Notiz oben): TP4 häufte Soldiers beim Archer an, um alle auf einmal zu treffen (B2); TP4 wurde beim Aufstehen erneut getroffen (B2).}
\end{itemize}

\paragraph{Interpretation}
Der Aussage stimmten vier von fünf Testpersonen eher oder völlig zu; der Kampf gegen Gruppen wurde als schwer, aber nicht unangenehm beschrieben. Unsere Umsetzung mit grossen und aggressiven Gruppen führte allerdings auch dazu, dass man nicht genau einschätzen konnte, wie viel Schaden eine Gruppe anrichten kann. Wird man in der Luft getroffen und muss wieder aufstehen, kann es passieren, dass man von einer Gruppe umzingelt wird und viel Schaden erleidet, was möglicherweise zu stark bestraft.

In unserer Umsetzung haben wir vor allem die Anzahl der gleichzeitig gespawnten Gegner begrenzt und weniger, wie viele davon aktiv angreifen. Bei unseren eigenen Tests vor dem Playtest fühlte sich das besser an, es könnte aber sein, dass wir damit übertrieben haben. Mit diesem Vorschlag wollten wir erreichen, dass sich der Kampf schwer anfühlt, die Schwierigkeit aber dennoch machbar bleibt. Schlussendlich haben wir es ähnlich wie \game{Ninja Gaiden} gelöst, wo die maximale Anzahl aktiver Gegner begrenzt wird (Abschnitt~\ref{sec:Vergleich}).

\begin{itemize}[leftmargin=1.5cm]
    \item \textcolor{red}{Urteil ergänzen, z.\,B.: Wirkung bestätigt (Kontrolle bleibt erhalten), aber über einen anderen Mechanismus als im Vorschlag formuliert. Das passt zum roten Hinweis in Abschnitt~\ref{sec:V2}.}
    \item \textcolor{red}{\enquote{zu stark bestraft}: Belege nennen (TP4 B2: beim Aufstehen erneut getroffen; TP5 Q13: Stunlock).}
\end{itemize}

\begin{itemize}[leftmargin=1.5cm]
    \item \textcolor{red}{\textbf{Neu -- Interpretation:}}
    \item \textcolor{red}{Urteil: für vier von fünf bestätigt; TP1 ist die Ausnahme. Kann man TP1s tiefen Wert mit dem Weglaufen (B2) und der Ecke (B6) erklären?}
    \item \textcolor{red}{Bezug Abschnitt~\ref{sec:V2}: Begrenzt wurde vor allem die Zahl der gleichzeitig lebenden Gegner (3--6), nicht die der Angreifer, weil der Gegnermanager aggressiv eingestellt ist.}
\end{itemize}

\subsubsection{V3: Ein taktisches Dilemma durch Nah- und Fernkämpfer}

\paragraph{Ziel}
Der Kampfabschnitt soll Nah- und Fernkämpfer kombinieren, sodass der Spieler entscheiden muss, welche Bedrohung er zuerst ausschaltet. Die Schwierigkeit soll dadurch taktisch statt rein mechanisch entstehen und schon kleine Gruppen anspruchsvoll machen (Abschnitt~\ref{sec:V3}).

\paragraph{Ergebnisse}
Im Interview gaben alle Testpersonen an, zuerst die Archers anzugreifen, wenn sich die Möglichkeit bietet. TP2 meinte, sie sei vor allem auf die Chain und Heavy Archers losgegangen und habe sich die Heavy Soldiers für den Schluss aufbewahrt, da die Archers mit höherer Wahrscheinlichkeit trafen und die anderen Gegner besser kontrollierbar waren. TP3 meinte, sie müsse nicht unbedingt auf die Soldiers losgehen, da diese einen ohnehin suchen und auf den Spieler zurennen. TP4 hat das taktische Dilemma zwar gesehen, meinte aber, es habe für sie keine allzu grosse Rolle gespielt; sie habe einfach, wie sie wollte, alles kurz und klein gehauen. Erst wenn sie merkte, dass es so nicht ging, und dazu gezwungen war, habe sie sich Gedanken zum Vorgehen gemacht.

Beobachtet wurde vor allem bei TP2 und TP3 das erwartete Verhalten: zuerst die Archers, dann die Soldiers, wobei die Soldiers den Weg zu den Archers blockierten. TP4 schlug dagegen zuerst auf die Soldiers ein und ging erst danach auf die Archers los. TP5 probierte Verschiedenes aus: Manchmal griff sie zuerst die Archers an, manchmal die Nahkämpfer. Ausschlaggebend waren für sie eher die Gegner auf den Boxen, weil sie sich so einen sicheren Platz zum Heilen verschaffen konnte.

\begin{itemize}[leftmargin=1.5cm]
    \item \textcolor{red}{Angepasst: Euer erster Satz lautete, das Dilemma sei von allen \emph{wahrgenommen und beobachtet} worden, dass alle zuerst auf die Archers gehen. Im Interview sagen das alle (I5), beobachtet wurde es aber nur bei TP2 und TP3 klar. Q3 als Tendenz ergänzen (TP4 neutral, die anderen stimmen zu).}
    \item \textcolor{red}{TP2 laut Anhang: zuerst Archers mit drei Pfeilen oder zielsuchenden Pfeilen; die \enquote{mit den Hämmern} zum Schluss, weil langsam und konterbar (I5). \enquote{Ignorierte} habe ich daher zu \enquote{für den Schluss aufbewahrt} geändert.}
\end{itemize}

\begin{itemize}[leftmargin=1.5cm]
    \item \textcolor{red}{\textbf{Neu -- Ergebnisse (Daten aller fünf Testpersonen):}}
    \item \textcolor{red}{\textbf{Q3:} TP2, TP3 = 5, TP1, TP5 = 4, TP4 = 3.}
    \item \textcolor{red}{\textbf{Interview (I5):} jetzt fünf von fünf nennen Archers zuerst. TP1: zuerst den Chain Archer, \enquote{weil ich den nicht mochte}, dann die Soldiers auf dem Weg, am Schluss die Heavy Archers, \enquote{Intuition}.}
    \item \textcolor{red}{\textbf{Beobachtung TP1 (B3):} wollte zu den Archers, wurde aber von den Gegnern dazwischen abgehalten, und griff zuerst die am höchsten platzierten Archers an. Das ist genau die Situation aus Abschnitt~\ref{sec:V3}; TP1 zu \enquote{beobachtet vor allem bei TP2 und TP3} hinzufügen.}
\end{itemize}

\paragraph{Interpretation}
Das taktische Dilemma wurde von allen Testpersonen wahrgenommen, und es gab verschiedene Ansätze, es zu lösen. Grundsätzlich liess sich beobachten, dass die Archers und das Räumen der Boxen Vorrang hatten. Aufgrund der unterschiedlichen Gegnertypen empfanden die einzelnen Testpersonen allerdings unterschiedliche Gegner als lästig.

Wir konnten aber auch beobachten, dass man das Spiel mit wenig Taktik durchspielen konnte. Zudem gab es widersprüchliche Antworten: Manche Testpersonen erkannten das taktische Dilemma zwar, handelten aber nicht immer danach. So war es oft ebenfalls zielführend, einfach den erstbesten Gegner anzugreifen, und der Schaden der Archers wurde teils ignoriert. In manchen Situationen wurde das Dilemma zwar erkannt, doch Schaden in Kauf zu nehmen war die mental einfachere Option und passte auch zum aggressiven Spielstil.

Umgesetzt wurde der Vorschlag ziemlich gut: Die Archers wurden oft manuell und häufig am Rand platziert. Da sie sich nicht bewegen, befand sich so oft ein Soldier zwischen Spieler und Archer. Ob der Archer dann zuerst angegriffen wurde, war unterschiedlich.

\begin{itemize}[leftmargin=1.5cm]
    \item \textcolor{red}{Urteil ergänzen, z.\,B.: wahrgenommen ja (Interview, 5 von 5), handlungsleitend nur teilweise (Beobachtung).}
    \item \textcolor{red}{\enquote{das Räumen der Boxen hatte Vorrang}: belegt ist das nur bei TP5 (B3). Präzisieren.}
    \item \textcolor{red}{Optional: Liegt das Dilemma eher an der Platzierung (V5) als an den Gegnerrollen? Bezug auf Frage D2 der Analysen (Abschnitt~\ref{sec:Vergleich}).}
\end{itemize}

\begin{itemize}[leftmargin=1.5cm]
    \item \textcolor{red}{\textbf{Neu -- Interpretation:}}
    \item \textcolor{red}{Urteil: wahrgenommen von allen (I5), handlungsleitend nur teilweise. TP1 stützt den Mechanismus aus Abschnitt~\ref{sec:V3} (Nahkämpfer versperren den Weg).}
\end{itemize}

\subsubsection{V4: Klar telegrafierte Angriffe}

\paragraph{Ziel}
Bevorstehende Gegnerangriffe sollen deutlich angekündigt werden. So soll der Spieler auch im Gruppenkampf rechtzeitig reagieren können und Treffer als eigene Fehler statt als Willkür erleben (Abschnitt~\ref{sec:V4}).

\paragraph{Ergebnisse}
Bei der Aussage, die Angriffe seien klar angekündigt worden (Q4), gehen die Meinungen am weitesten auseinander. Interessanterweise stimmte TP3, die sehr aggressiv und ohne viel nachzudenken spielte, völlig zu. TP2, die ebenfalls sehr aggressiv spielte und meistens erst reagierte, nachdem ein Treffer gelandet war, bewertete die Aussage neutral. Auch TP5, die eher taktisch spielte, bewertete die Aussage neutral, da sie die Ankündigungsanimationen der Archers verwirrend fand. TP4, die Projektilen auswich, wenn sie konnte, sonst aber alle Treffer in Kauf nahm, stimmte überhaupt nicht zu. Dazu meinte sie, die Archers seien besonders, da sie ausserhalb des Blickfelds angreifen könnten und es keinen Sound gebe, der den Angriff ankündigt. TP5 meinte zudem, die Tode seien wegen der Ankündigungen nicht unfair gewesen: Man habe kommen sehen, dass man stirbt, und es habe sich eher wie ein selbstverschuldetes Problem angefühlt.

\begin{itemize}[leftmargin=1.5cm]
    \item \textcolor{red}{\enquote{ausserhalb des Blickfelds}: Im Anhang steht bei TP4 nur \enquote{ohne wirkliche Prediction oder Sound-Effekte} (I5). Falls das aus der Aufnahme stammt, im Anhang ergänzen.}
    \item \textcolor{red}{Ergänzen: TP3 sah die Pfeile kommen, konnte wegen der Combos aber nicht ausweichen (I6, Q13).}
\end{itemize}

\begin{itemize}[leftmargin=1.5cm]
    \item \textcolor{red}{\textbf{Neu -- Ergebnisse (Daten aller fünf Testpersonen):}}
    \item \textcolor{red}{\textbf{Q4:} TP3 = 5, TP2, TP5 = 3, \textbf{TP1 = 2}, TP4 = 1. Die Streuung bleibt die grösste.}
    \item \textcolor{red}{\textbf{TP1 ergänzen:} versuchte auszuweichen, ignorierte Angriffe aber oft (B4); wurde beim Springen auf Boxen von Archers abgeschossen (B4, I1b); empfand Treffer am Anfang als unfair, \enquote{bis ich gecheckt habe, was die neuen Gegner machen} (I6).}
    \item \textcolor{red}{TP5 fand die Ankündigungsanimation der Archers verwirrend, weil sie aussieht, als würden sie laufen (B4).}
\end{itemize}

\paragraph{Interpretation}
Bei diesem Vorschlag gab es die grösste Streuung in den Antworten. Dass TP3 völlig zustimmte, kann damit zusammenhängen, dass ihr die meisten Treffer egal waren und ihr die Ankündigungen deshalb genügten. Dass TP2 erst nach einem Treffer reagierte und nur neutral bewertete, könnte darauf hindeuten, dass die Ankündigungszeit etwas zu kurz war. Dass TP4 überhaupt nicht zustimmte, kann damit zusammenhängen, dass sie die Information nicht brauchte und deshalb nicht darauf achtete, oder auch daran liegen, dass die Combos zu träge sind und sie deshalb auf die Ankündigungen nicht reagieren konnte. TP5 meinte, die Ankündigungen hätten die Tode nicht unfair, sondern weniger schlimm gemacht, da man sie kommen sah und sie nicht aus dem Nichts kamen.

Die Streuung hängt wahrscheinlich mit den Spielstilen zusammen: Taktische Testpersonen bewerteten die Ankündigungen eher tief und wünschten sich mehr davon, während aggressive Testpersonen, die sich wenig für den erlittenen Schaden interessierten, sie als ausreichend empfanden. Die Angriffe der Soldiers waren gut genug zu erkennen, da diese zum Spieler laufen müssen, um Schaden zu verursachen. Das galt nicht für die Archers, die auch von hinten ohne Ankündigung angreifen konnten.

Dabei müssen zwei Probleme getrennt werden: zum einen die Wahrnehmung, also die Ankündigung an sich, und zum anderen die Reaktionsmöglichkeit, also ob man auf die Ankündigung überhaupt reagieren kann.

\begin{itemize}[leftmargin=1.5cm]
    \item \textcolor{red}{\textbf{Widerspruch zu den Daten:} \enquote{aggressive Testpersonen fanden es genug} trifft nur auf TP3 zu (Q4 = 5). TP4 spielte ebenfalls eher aggressiv (B1, I4: \enquote{kopfvoran}) und bewertete am tiefsten (1); TP2 spielte aggressiv und bewertete neutral (3). Die Spielstil-Erklärung so nicht halten oder deutlich einschränken.}
    \item \textcolor{red}{Urteil ergänzen. Und: Liegt es am Vorschlag oder an der Umsetzung (nur Animation, kein Sound, keine Partikel)? Bezug auf das Telegrafieren über mehrere Kanäle in \game{NieR: Automata} (Frage M2) und auf eure Begründung in Abschnitt~\ref{sec:V4}.}
    \item \textcolor{red}{Den Satz zu TP5 nicht wiederholen, wenn dieser schon bei den Ergebnissen steht; hier nur interpretieren.}
\end{itemize}

\begin{itemize}[leftmargin=1.5cm]
    \item \textcolor{red}{\textbf{Neu -- Interpretation:}}
    \item \textcolor{red}{Urteil fehlt noch. TP1 (Q4 = 2) ist ein weiteres Gegenbeispiel zu \enquote{aggressive Testpersonen fanden die Ankündigungen ausreichend}: Mit TP1 bewerteten zwei aggressiv spielende Testpersonen tief (TP1, TP4), eine hoch (TP3).}
    \item \textcolor{red}{Bezug Abschnitt~\ref{sec:V4}: Dort argumentiert ihr, die Animationen reichten, weil der Spieler ohnehin ständig mit einem Angriff rechnen muss. Stützen die Daten das?}
\end{itemize}

\subsubsection{V5: Lesbare Platzierung, dynamisches Verhalten}

\paragraph{Ziel}
Die Gegner sollen zu Beginn lesbar und teils symmetrisch platziert sein, sich danach aber dynamisch verhalten, etwa den Spieler umzingeln. Die Anfangsaufstellung soll Überblick geben, das anschliessende Verhalten den Kampf unter Druck und unvorhersehbar halten (Abschnitt~\ref{sec:V5}).

\paragraph{Ergebnisse}
Hier stimmten TP1, TP2 und TP3 eher oder völlig zu, TP4 und TP5 eher nicht (jeweils Q5). TP2 kletterte zwischen den Wellen auf Boxen, um den Raum zu überblicken. TP3 griff einfach an, was vor ihr stand; die vorhandene Information reichte ihr also aus. In Raum~10, einer engen, verwinkelten Gasse, fand TP3 als einzige Testperson den versteckten Archer eine Zeit lang nicht. Auch TP5 griff eher an, was vor ihr lag, und verstand das dynamische Verhalten der Gegner. Sie meinte, man wisse nicht, wo die Gegner spawnen, bevor man den Raum betreten habe. TP5 kreiste zudem hin und wieder eine Weile im Raum, bevor sie angriff, um die Gegner zu beobachten, da sie so nicht getroffen werden konnte; auf diese Weise sammelte sie die Gegner auch an. TP4 fand es unangenehm, von hinten, also ausserhalb des Sichtfelds, getroffen zu werden. Sie schaute sich eine Welle eher lange an, bevor sie angriff.

TP3 bemerkte, dass Gegner ausserhalb des Sichtfelds spawnen, meinte aber, man bemerke sie früh genug, um reagieren zu können. Sie stellte auch fest, dass die Archers eher erhöht spawnen und dass es definitiv eine manuell platzierte Konstante geben müsse; es seien also deutliche Muster erkennbar. TP2 gefiel die Verteilung der Gegner im Raum am besten. TP5 meinte nach einem Tod, es fühle sich richtig gut an, wenn man wisse, wo ein Gegner spawnen wird, und ihn direkt zu Beginn ausschalten könne.

\begin{itemize}[leftmargin=1.5cm]
    \item \textcolor{red}{Ergänzen: TP4 starb in Raum~12, weil sie einen Archer übersehen hatte (B6). Optional: TP5 wurde von der Gegnermasse fast in einen Türrahmen gedrängt (B6).}
    \item \textcolor{red}{Verweise auf den Anhang ergänzen: TP2 (B6, Q12), TP3 (B6, I8), TP4 (B6, B11), TP5 (B2, B6, Q12).}
    \item \textcolor{red}{\enquote{Raum~10}: Im Anhang steht bei TP3 nur \enquote{versteckter Heavy Archer} ohne Raumnummer. Falls aus der Aufnahme, dort ergänzen.}
\end{itemize}

\begin{itemize}[leftmargin=1.5cm]
    \item \textcolor{red}{\textbf{Neu -- Ergebnisse (Daten aller fünf Testpersonen):}}
    \item \textcolor{red}{\textbf{TP1 ergänzen (B6):} sagte die Platzierung der Gegner auch in neuen Räumen voraus; sprang bei einer neuen Welle zuerst heraus und nahm Abstand; griff an, was vor ihr lag.}
\end{itemize}

\paragraph{Interpretation}
Auch hier gehen die Meinungen auseinander. TP2 und TP3 stimmten zu: TP3 erkannte Muster in den Wellen, und TP2 gefiel die Verteilung der Gegner in den Räumen. TP4 und TP5 fanden dagegen, dass zwar das dynamische Verhalten der Gegner sichtbar war, die Anfangsposition aber nicht unbedingt. Manchmal wurden auch Archers übersehen. Unerwartete Treffer von hinten gab es trotzdem, und die Formation der Gegner lässt sich nicht sehen, bevor man den Raum betritt. Die Antworten sind zudem schwer zu deuten, da die Aussage Q5 zwei Dinge gleichzeitig abfragt: die Überschaubarkeit zu Beginn und das danach unvorhersehbare, dynamische Verhalten.

\begin{itemize}[leftmargin=1.5cm]
    \item \textcolor{red}{Urteil ergänzen.}
    \item \textcolor{red}{Bezug zur Umsetzung: Weil die Gegner erst nach dem Betreten spawnen, liegt euer Prototyp näher an \game{Ninja Gaiden} (Spawn ausser Sicht, Frage A3) als an der sichtbaren, symmetrischen Aufstellung aus \game{NieR: Automata}. Ein Satz dazu.}
    \item \textcolor{red}{Die Doppelfrage Q5 auch in den Limitationen (Kapitel~\ref{sec:Schlussfolgerungen}) erwähnen.}
    \item \textcolor{red}{Der erste Absatz wiederholt die Ergebnisse direkt darüber; kürzen.}
\end{itemize}

\begin{itemize}[leftmargin=1.5cm]
    \item \textcolor{red}{\textbf{Neu -- Interpretation:}}
    \item \textcolor{red}{Urteil fehlt noch. TP1 sagte Platzierungen voraus (B6), ebenso TP3 (I8) und TP4 (B11): Die Platzierung war also erkennbar, sobald man die Muster kannte. Gleichzeitig der Überraschungseffekt beim Betreten, den ihr in Abschnitt~\ref{sec:V5} beschreibt (TP1 nimmt bei neuen Wellen erst Abstand).}
\end{itemize}

\subsubsection{V6: Der Raum als Druckmittel}

\paragraph{Ziel}
Die Grösse des Kampfraums soll gezielt zwischen Enge und Weite wechseln. Ein enger Raum soll den Spieler zu methodischerem Spiel zwingen, während ihm ein weiter Raum taktische Distanz erlaubt (Abschnitt~\ref{sec:V6}).

\paragraph{Ergebnisse}
Hier gehen die Meinungen etwas auseinander. TP2 spielte in grösseren und kleineren Räumen unterschiedlich: In grossen Räumen spielte sie aggressiver, in kleineren weniger aggressiv und eher taktisch, da es dort wenig Platz zum Heilen gab und das Leben dadurch wertvoller war. Bei TP3 war es unterschiedlich: In manchen kleinen Räumen spielte sie genau gleich aggressiv, in anderen hektischer. TP5 bemerkte dagegen einen deutlichen Unterschied in kleinen Räumen. Wenn diese verwinkelt waren, empfand TP5 sie als gruselig und schaute deshalb vorsichtig um die Ecken.

TP2 meinte, es gebe eine gute Mischung der Raumgrössen. Die kleinen Räume seien tendenziell schwerer, weil es keine Zufluchtsmöglichkeit gebe; ihr Lieblingsraum sei aber Raum~10 mit der Brücke gewesen. TP3 meinte, kleine Räume seien einfacher zu überblicken gewesen, man habe direkt gewusst, was gespawnt sei und was auf einen zukomme. In grösseren Räumen mit Boxen und Brücke habe es dagegen mehr Möglichkeiten gegeben, sich im Raum zu bewegen. Für TP4 spielte die Raumgrösse keine Rolle, weil sie einfach alles kopfvoran versucht hat. TP5 meinte, kleine Räume seien übersichtlicher, man könne aber in Ecken gedrängt werden; die grossen seien schwerer gewesen, weil es dort mehr Gegner gab. Wie bereits bei V2 beschrieben, hatte TP5 etwas Mühe mit grösseren Gegnermengen.

\begin{itemize}[leftmargin=1.5cm]
    \item \textcolor{red}{Korrigiert: \enquote{für TP2 spielte es keine Rolle} war laut Anhang TP4 (I4), bitte prüfen.}
    \item \textcolor{red}{Q6 als Tendenz ergänzen: Die Antworten gehen weit auseinander (TP2 und TP5 stimmen völlig zu, TP3 neutral, TP4 eher nicht).}
    \item \textcolor{red}{\enquote{gruselig} (TP5) steht nicht im Anhang -- falls aus der Aufnahme, dort ergänzen.}
\end{itemize}

\begin{itemize}[leftmargin=1.5cm]
    \item \textcolor{red}{\textbf{Neu -- Ergebnisse (Daten aller fünf Testpersonen):}}
    \item \textcolor{red}{\textbf{Q6:} TP2, TP5 = 5, TP1 = 4, TP3 = 3, TP4 = 2. Die Antworten gehen weit auseinander.}
    \item \textcolor{red}{\textbf{TP1 ergänzen (I4):} Im Brückenraum gab es mehr Möglichkeiten, sich zu verstecken und auszuweichen, stressiger war er aber nicht, die Gegner waren \enquote{verteilter}. Beim Beobachtungskriterium B7 gibt es bei TP1 keinen Eintrag.}
    \item \textcolor{red}{TP5 (I4): kleine Räume übersichtlicher, aber Gefahr, in eine Ecke gedrängt zu werden; grosse Räume schwerer, weil mehr Gegner. TP5 war in Raum~10 vorsichtig um die Ecken (B7).}
\end{itemize}

\paragraph{Interpretation}
Dass TP3 in kleinen Räumen teils gleich aggressiv, teils hektischer spielte, liegt wahrscheinlich daran, dass sie wegen ihres aggressiven Spielstils ohnehin wenig Platz brauchte und den zusätzlichen Platz der Prospect Spaces nur in Momenten mit wenig Leben oder zu grosser Gefahr vermisste.

\begin{itemize}[leftmargin=1.5cm]
    \item \textcolor{red}{Noch offen: Urteil, Erklärung und Bezug. Anregungen: Q6 streut stark; TP3 und TP5 fanden kleine Räume \emph{übersichtlicher}; der Druck in engen Räumen entstand eher durch fehlenden Platz zum Heilen (TP2) als durch die Nahkämpfer. Passt das zu eurer Begründung in Abschnitt~\ref{sec:V6} und zu Totten?}
\end{itemize}

\begin{itemize}[leftmargin=1.5cm]
    \item \textcolor{red}{\textbf{Neu -- Interpretation:}}
    \item \textcolor{red}{Urteil, Erklärung und Bezug fehlen noch. TP1 und TP5 beschreiben eher den Unterschied in Möglichkeiten und Übersicht als einen Unterschied im Druck. Passt das zu eurer Einteilung nach Bewegungsfreiheit in Abschnitt~\ref{sec:V6}?}
\end{itemize}

\subsubsection{V7: Wellenstruktur mit Erholungsphasen}

\paragraph{Ziel}
Der Kampfabschnitt soll in Wellen mit kontrollierter Eskalation und kurzen Erholungsphasen dazwischen aufgebaut sein. So sollen sich Anspannung und Entlastung abwechseln, und die Schwierigkeit soll planbar ansteigen (Abschnitt~\ref{sec:V7}).

\paragraph{Ergebnisse}
Alle fünf Testpersonen organisierten ihre Erholungszeit selbst. TP2 machte Pause, bevor sie in den nächsten Raum ging, und nutzte die Boxen. TP3 ist gestorben, weil sie sich anfangs nicht geheilt hat; nachdem sie das bemerkt hatte, legte sie am Ende eines Raumes eine Ruhephase zum Heilen ein. Auch TP4 heilte sich vor dem Raumwechsel. TP5 ging dabei sogar etwas taktischer vor und liess oft einen Gegner am Ende einer Welle am Leben, um sich vor der nächsten Welle zu heilen, da man nicht genau wusste, was danach kommt.

TP2 meinte, es sei schwer, während eines Encounters Heilpausen zu finden. TP3 meinte, es gebe eine gute Balance; sobald alle Fernkämpfer besiegt seien, werde es entspannt. TP4 empfand den Verlauf der Kampfabschnitte ebenfalls als angenehm: Sie sei immer auf das Spiel konzentriert gewesen, und nichts habe ihre Vertiefung gestört. TP5 meinte, die Spannung hänge vom eigenen Leben ab, da man mit wenig Leben angespannter sei.

\begin{itemize}[leftmargin=1.5cm]
    \item \textcolor{red}{TP5 \enquote{liess oft einen Gegner am Leben}: Im Anhang steht \enquote{heilt sich, wenn am Ende einer Welle nur noch wenige Gegner übrig sind} (B8). Prüfen und angleichen.}
    \item \textcolor{red}{Q7 als Tendenz ergänzen (TP3 und TP5 stimmen völlig zu, TP2 und TP4 neutral). TP2 sprach zudem von \enquote{überlappenden Wellen} (I7) -- wichtig für die Diskussion.}
\end{itemize}

\begin{itemize}[leftmargin=1.5cm]
    \item \textcolor{red}{\textbf{Neu -- Ergebnisse (Daten aller fünf Testpersonen):}}
    \item \textcolor{red}{\textbf{Q7:} TP1, TP3, TP5 = 5, TP2, TP4 = 3.}
    \item \textcolor{red}{\textbf{TP1 ergänzen (B8):} ging auf die Boxen, um Abstand zu nehmen, und heilte hinter Boxen in Deckung; heilte regelmässig, sobald die Energieleiste voll war; vergass manchmal, vor dem Raumwechsel zu heilen, und heilte dann direkt im nächsten Raum.}
    \item \textcolor{red}{Erholungsmittel aus Abschnitt~\ref{sec:V7} im Vergleich: Boxen genutzt von TP1, TP2 und TP5 (B8); \enquote{letzten Gegner leben lassen} so nicht beobachtet, am nächsten kommt TP5, die heilt, wenn nur noch wenige Gegner übrig sind (B8).}
    \item \textcolor{red}{Interview I7: TP1 \enquote{gegen Ende immer cooler}; TP2 \enquote{überlappende Wellen}; TP3 gute Balance; TP4 \enquote{immer locked in}; TP5 hängt von den Lebenspunkten ab.}
\end{itemize}

\paragraph{Interpretation}
Dass TP2 es schwer fand, Heilpausen zu finden, liegt daran, dass sie die Wellenstruktur des Spiels nicht erkannt hat und deshalb nicht davon profitierte. \textcolor{red}{[Belegt? Hinweis: TP2 bemerkte das Ende eines Encounters nicht (B10). Sonst vorsichtiger formulieren: \enquote{könnte daran liegen}.]}

\begin{itemize}[leftmargin=1.5cm]
    \item \textcolor{red}{Noch offen: Urteil, Erklärung und Bezug. Anregungen: Alle organisierten ihre Erholung selbst, vor allem zwischen den Räumen; zwischen den Wellen gibt es in der Umsetzung kaum Pausen (Abschnitt~\ref{sec:V7}); fehlende Signale für das Ende einer Welle bzw.\ eines Raums (Abschnitt~\ref{sec:ErgAllgemein}). Vergleich: kurze Entlastung am Ende jeder Welle in \game{Ninja Gaiden} (Frage D1).}
\end{itemize}

\begin{itemize}[leftmargin=1.5cm]
    \item \textcolor{red}{\textbf{Neu -- Interpretation:}}
    \item \textcolor{red}{Urteil: Erholung haben sich alle selbst organisiert, meist zwischen den Räumen und mit Boxen, so wie in Abschnitt~\ref{sec:V7} vorgesehen. Das Lebenlassen des letzten Gegners wurde kaum genutzt.}
\end{itemize}

\subsubsection{V8: Gestaffelte Einführung neuer Gegner}

\paragraph{Ziel}
Jeder neue Gegnertyp soll zuerst isoliert auftreten und erst danach mit anderen kombiniert werden. So soll der Spieler einen Gegner in Ruhe verstehen können, bevor die Kombination ihn fordert (Abschnitt~\ref{sec:V8}).

\paragraph{Ergebnisse}
Hier stimmten alle Testpersonen völlig zu (Q8). Alle erkannten neue Gegnertypen bei ihrer Einführung an der Farbe, am Aussehen oder an ihren Angriffen. Wie sie mit ihnen umgingen, war allerdings unterschiedlich: TP2 hielt meist kurz inne und analysierte, ob sich noch etwas im Raum befand. TP4 beobachtete den Chain Archer zuerst, ohne anzugreifen, um zu sehen, was er macht. TP3 und TP5 griffen neue Gegner sofort an. TP3 dachte bei der Einführung des Archers, es handle sich einfach um eine stärkere Version der Soldiers, und wusste anhand des Aussehens anfangs nicht, was die neuen Gegner tun.

\begin{itemize}[leftmargin=1.5cm]
    \item \textcolor{red}{\textbf{Neu -- Ergebnisse (Daten aller fünf Testpersonen):}}
    \item \textcolor{red}{\textbf{Q8:} alle fünf = 5.}
    \item \textcolor{red}{\textbf{TP1 ergänzen (B9):} griff die ersten Archers nicht an, sondern wich aus; erkannte neue Gegner von Weitem an der Farbe, den Chain Archer an seinen Angriffen; bemerkte sofort, dass der Heavy Soldier ungenau trifft und der Heavy Archer weniger Leben hat.}
    \item \textcolor{red}{\textbf{Kombinieren (Braiding):} TP1 sagte, gegen Ende sei es \enquote{immer cooler} geworden, weil schon verschiedene Gegner eingeführt waren und es verschiedene Kombinationen gab (I7). Das ist der einzige direkte Beleg für den zweiten Teil des Vorschlags.}
    \item \textcolor{red}{Gegenstimme: TP3 erkannte neue Gegner, wusste aber nicht genau, was sie machen (I1b).}
\end{itemize}

\paragraph{Interpretation}
\begin{itemize}[leftmargin=1.5cm]
    \item \textcolor{red}{Noch zu schreiben. Anregungen: einstimmige Zustimmung (Q8), die stärkste im ganzen Fragebogen; alle erkennen neue Gegner. Aber: Erkennen heisst nicht verstehen (TP3 I1), und Q8 misst, ob die Reihenfolge bemerkt wurde, nicht, ob sie beim Verstehen geholfen hat. Geprüft wurde vor allem die Einführung, weniger das spätere Kombinieren (Braiding, Abschnitt~\ref{sec:Scaffolding}).}
\end{itemize}

\begin{itemize}[leftmargin=1.5cm]
    \item \textcolor{red}{\textbf{Neu -- Interpretation:}}
    \item \textcolor{red}{Urteil: Einführung klar bestätigt (fünf von fünf). Fürs Kombinieren gibt es nur einen Beleg (TP1 I7).}
\end{itemize}

\subsection{Weitere für die Erkenntnisfrage relevante Beobachtungen} \label{sec:WeitereBeobachtungen}

\paragraph{Ergebnisse}
Weitere Beobachtungen betreffen die Orientierung, die nicht immer reibungslos verlief: Oft wurde nicht bemerkt, dass ein Encounter beendet war, und offene Türen wurden übersehen, da es keinen audiovisuellen Hinweis darauf gab. Bei der Steuerung war den Testpersonen die Tiefe der Eingaben nicht immer intuitiv. So kann man die Dash-Taste gedrückt halten; Testpersonen, die das aus anderen Spielen kannten, waren davon begeistert, während die anderen die Taste eher spammten. Beim Springen verhielt es sich ähnlich. Auch die Reichweite der eigenen Waffe einzuschätzen, fiel etwas schwer, und die Dauer der Heilung war nicht ganz nachvollziehbar. TP2 schlug vor, statt der mündlichen Erklärung ein Tutorial einzubauen. TP4 und TP5 versuchten, mit Objekten zu interagieren, und suchten beide nach Glitches.

Die Stimmung war in allen Playtests positiv: TP3 und TP4 lachten viel, und TP3 und TP5 spielten nach dem Ende einfach weiter. TP2 war hin und wieder genervt, wenn sie beim Heilen unterbrochen wurde. Fast jede Testperson hatte einen Gegnertyp, der sie besonders störte. TP5 meinte, sie sei nie unter Stress gewesen.

Gewissermassen gab es zwei Bewegungsmodi, die normale Fortbewegung und das Kämpfen. Der Wechsel zwischen diesen Modi bereitete den Testpersonen hin und wieder Schwierigkeiten, und die Reichweite des Spielers war etwas schwer einzuschätzen.

\begin{itemize}[leftmargin=1.5cm]
    \item \textcolor{red}{\textbf{Faktencheck mit dem Anhang:}}
    \item \textcolor{red}{Glitches suchten laut Anhang TP3 (I7: \enquote{Bugs zu finden hat mir Freude bereitet}) und TP5 (B13), nicht TP4. TP4 und TP5 wollten mit Objekten interagieren (B13).}
    \item \textcolor{red}{Nicht im Anhang belegt: TP4 lachte viel, TP5 spielte nach dem Ende weiter, Testpersonen änderten ihre Meinung zu Block/Dodge, Begeisterung über das Gedrückthalten von Dash. Falls aus den Aufnahmen, im Anhang ergänzen.}
    \item \textcolor{red}{Pauschale Aussagen (\enquote{die Testpersonen}, \enquote{oft}) durch \enquote{x von 5} bzw.\ TP-Nummern mit Verweis ersetzen, z.\,B. offene Türen übersehen: TP2 und TP4 (B10), TP5 nicht (B14); Springen: TP2 (Q13); Übergang Bewegung/Kampf: TP5 (B10); Reichweite: TP3 (Q13), TP5 (B10); Heilungsdauer: TP4 (B10); Boss erwartet: TP2 (Q14), TP3 (I2).}
\end{itemize}

\paragraph{Interpretation}
\begin{itemize}[leftmargin=1.5cm]
    \item \textcolor{red}{Hier gehört hin, was keinem Designvorschlag zugeordnet ist, das Spielerlebnis aber geprägt hat. Kandidaten aus den Ergebnissen oben: Trägheit der Combos und fehlender Abbruch, Treffer in der Luft und Aufstehen im Stunlock, übersehene Türen bzw.\ fehlendes Signal am Ende eines Encounters, Übergang zwischen Bewegung und Kampf, Reichweite, Heilungsdauer. Jeweils: was beobachtet, wie viele Testpersonen, warum relevant für Spielfluss, Schwierigkeit oder Spielgefühl.}
\end{itemize}

\subsection{Gesamteindruck: Spielgefühl, Spielfluss und Schwierigkeit} \label{sec:Gesamteindruck} \label{sec:ErgAllgemein}

\paragraph{Ergebnisse}
Der Gesamteindruck der Testpersonen in den drei Kategorien fällt ziemlich ähnlich aus. Das Spielgefühl beurteilten alle als eher befriedigend (Q9). Den Spielfluss empfanden ebenfalls alle als zusammenhängend (Q10); dabei meinten sie, dass die Räume zwar jeweils für sich stehen und wenig miteinander verbindet, der Spielfluss insgesamt aber gepasst habe. Auch die Schwierigkeit fanden alle angemessen (Q11): Alle sind ein- bis zweimal gestorben, aber nicht wiederholt an derselben Stelle, und es gab viele Momente, in denen sie knapp überlebten.

Als das, was sich am besten anfühlte (Q12), nannten die Testpersonen die Gegnerverteilung und die Spawns, etwa wenn man Gegner direkt beim Spawnen besiegen kann. Auch die taktischen Entscheidungen, das Heilen und das Weglaufen wurden als interessant genannt. Horden zu vernichten, was der Kern des Spiels ist, habe Spass gemacht und sich gut angefühlt. Positiv erwähnt wurde zudem die Tiefe durch die unterschiedlichen Combos, also Light und Heavy mit ihren verschiedenen Schadenswerten.

Als unfair oder unangenehm empfanden die Testpersonen (Q13), dass man Combos nicht abbrechen und deshalb nicht so einfach auf Gegner reagieren konnte. Auch in einer Horde von Gegnern eingeklemmt zu sein, fühlte sich unbefriedigend an. Eine Testperson meinte zudem, das Springen fühle sich nicht wie erwartet an; das deckt sich mit einer anderen Aussage, wonach der Übergang zwischen Bewegung und Kampf etwas unstimmig sei.

Im Vergleich mit ihren Erwartungen (Q14) waren die Testpersonen vom Umfang und von der Gegnervielfalt positiv überrascht. Erwartet wurde ein Boss, der nicht kam; zudem wurden hin und wieder Moves wie Block oder Dodge gewünscht. Es gab aber auch Stimmen, die ihre Meinung am Ende änderten und meinten, dass diese Moves nicht ins Spiel gepasst hätten. Ausserdem wurde erwartet, dass man als Spieler noch stärker sei.

\begin{itemize}[leftmargin=1.5cm]
    \item \textcolor{red}{\textbf{Faktencheck mit dem Anhang:}}
    \item \textcolor{red}{\enquote{Alle sind ein- bis zweimal gestorben}: TP2 ist laut Interview gar nicht gestorben (I3), TP1 einmal (Q12), TP5 dreimal (I3), TP4 mindestens zweimal (I1). Anpassen.}
    \item \textcolor{red}{\enquote{Den Spielfluss empfanden alle als zusammenhängend}: TP3 gab \enquote{1 von Raum zu Raum, 4 im Kampf} an (Q10). Das deckt sich mit eurem Satz über die getrennten Räume -- so formulieren, statt \enquote{alle}.}
\end{itemize}

\paragraph{Interpretation}
\begin{itemize}[leftmargin=1.5cm]
    \item \textcolor{red}{Gliederung nach den drei Qualitäten aus der Erkenntnisfrage, je ein kurzer Absatz: Wurde die Qualität erreicht, und welche Vorschläge haben dazu beigetragen oder gestört?}
    \item \textcolor{red}{\textbf{Spielfluss:} Q10, B14; Beitrag u.\,a. von V2, V5, V7, V8.}
    \item \textcolor{red}{\textbf{Schwierigkeit:} Q11; Beitrag u.\,a. von V2, V3, V4, V6, V8.}
    \item \textcolor{red}{\textbf{Spielgefühl:} Q9, Q12--Q14; Beitrag u.\,a. von V1, V4, V5.}
    \item \textcolor{red}{Die Zuordnung stammt aus dem \enquote{Der Nutzen betrifft \ldots} in Kapitel~\ref{sec:Umsetzung}. Dieser Abschnitt ist die Brücke zur Beantwortung der Erkenntnisfrage in Abschnitt~\ref{sec:Beantwortung}.}
\end{itemize}

\end{document}
